\chapter{Introduction}

Banks and financial institutions process hundreds of millions of transactions every single day. Hidden inside this enormous volume are fraudsters, money launderers, identity thieves, and organised criminal networks, all attempting to move illegitimate funds or steal from genuine customers. The scale of the problem is significant: banking fraud is estimated to cost more than \$45 billion annually worldwide, cross-border payment volumes are projected to exceed \$200 trillion by 2025, and regulators impose roughly \$5 billion in fines every year on banks that fail to meet Anti-Money Laundering (AML) obligations. In the United States alone, consumers reported losing \$12.5 billion to fraud in 2024.

Despite this scale, most banks continue to rely on \textit{rule-based} fraud detection systems, in which engineers manually author static rules such as ``if a transaction exceeds a fixed amount \textit{and} the account is newly opened \textit{and} the recipient is overseas, flag it for review.'' While simple to implement, this approach is fundamentally reactive: it only catches fraud after a transaction has already been attempted, it cannot adapt to new criminal tactics without manual rule updates, and it treats each transaction in isolation rather than as part of a broader behavioural or relational pattern. In practice, over 95\% of the alerts generated by such systems turn out to be false positives, meaning analysts spend the majority of their time investigating entirely legitimate activity while genuine fraud can slip through amid the noise.

A second, equally important limitation of existing systems is data fragmentation. A bank typically holds transaction histories, Know-Your-Customer (KYC) records, login and device behaviour, geolocation data, and watchlist information, but these data sources are usually managed by separate teams using separate systems that rarely communicate with one another. Fraud investigators look at transactions, compliance teams look at KYC documentation, and security teams look at login and device data, but no single view brings all of this together. This is precisely where sophisticated fraud rings and money-laundering networks hide: an individual account may look completely normal in isolation, while sharing a phone number, device fingerprint, or IP address with several other accounts reveals a coordinated network of criminal activity.

This project proposes an AI-powered platform that addresses both limitations directly. Instead of evaluating a transaction in isolation against static thresholds, the platform unifies transaction, KYC, and behavioural data into a single knowledge graph, allowing hidden relationships between seemingly unrelated accounts to be surfaced. On top of this unified view, an AI-based risk-scoring model evaluates each transaction in real time using both its own features and the surrounding network context, producing an explainable risk score rather than a single-rule pass/fail decision. Finally, once a transaction or account is flagged, a Retrieval-Augmented Generation (RAG) based compliance assistant automatically drafts a Suspicious Activity Report (SAR) grounded in the actual regulatory text relevant to the case, reducing a task that otherwise takes an analyst two to four hours down to a matter of minutes.

The central shift this project makes is moving away from asking ``is this transaction, on its own, suspicious?'' towards asking ``does this account, in the context of its entire network of relationships and behavioural history, match a known criminal pattern?'' This entity-level, network-aware framing, combined with a real-time, explainable, and RAG-grounded compliance pipeline, forms the basis of the design discussed throughout the remainder of this report.
